\documentclass[10pt]{article}
\usepackage[letterpaper,margin=0.58in]{geometry}
\usepackage{amsmath,amssymb,amsthm,mathtools,mathrsfs}
\usepackage{enumitem}
\usepackage{microtype}
\usepackage[colorlinks=true,linkcolor=blue,urlcolor=blue]{hyperref}

\setlist{nosep,leftmargin=*}
\setlength{\parindent}{0pt}
\setlength{\parskip}{3pt}
\setcounter{secnumdepth}{1}
\linespread{0.96}

\newtheorem*{theorem}{Theorem}
\newcommand{\g}{\mathfrak g}
\newcommand{\uP}{\mathfrak u}
\newcommand{\OO}{\mathcal O}
\newcommand{\K}{K_{X_0}}
\newcommand{\A}{\mathscr A}
\newcommand{\Amr}{\mathscr A^{\mathrm{mr}}}
\newcommand{\rad}{\operatorname{rad}}
\newcommand{\im}{\operatorname{Im}}
\newcommand{\Hom}{\operatorname{Hom}}

\title{\vspace{-2.2em}Short-Kappa Maximal-Parabolic Induction\\
\large Dehydrated Proof}
\author{Condensed from the Rethlas-verified blueprint}
\date{May 4, 2026}

\begin{document}
\maketitle
\vspace{-1.5em}
\small

Let \(F=\mathbb R\) and \(\epsilon=\pm1\).  Let
\((V_0,\langle\,,\,\rangle_0)\) be non-degenerate \(\epsilon\)-symmetric,
let \(G_0=\operatorname{Isom}(V_0)\), and let \(X_0\in\mathfrak g_0\) be
nilpotent.  Put
\[
c_1=\dim\ker X_0,\qquad
c_2=\dim(\ker X_0\cap\im X_0)=\dim(\ker X_0^2/\ker X_0).
\]
Let \(E,E'\) be \(\kappa\)-dimensional real vector spaces with perfect pairing
\(\lambda:E\times E'\to\mathbb R\), and assume
\[
c_2\le \kappa<c_1,\qquad d:=c_1-\kappa.
\]
On \(V=E\oplus V_0\oplus E'\) use
\[
\langle e+v+e',f+w+f'\rangle
=\lambda(e,f')+\epsilon\lambda(f,e')+\langle v,w\rangle_0 .
\]
Let \(G=\operatorname{Isom}(V)\), let \(P\subset G\) stabilize \(E\), and let
\[
\uP=\{Y\in\g:Y(E)=0,\ Y(E^\perp)\subset E,\ Y(V)\subset E^\perp\}.
\]
We compute
\[
\operatorname{Ind}_P^G(\OO_0)=\overline{G\cdot(\OO_0+\uP)},
\qquad \OO_0=G_0\cdot X_0 ,
\]
with real-topological closure.  Since \(X_0\) is skew-adjoint,
\[
(\im X_0)^\perp=\ker X_0.
\]
Thus the form on \(\ker X_0\) has radical
\(R:=\ker X_0\cap\im X_0\), and descends to a non-degenerate
\(\epsilon\)-symmetric form on
\[
\K:=\ker X_0/R .
\]
Let \(\A=\operatorname{Gr}_d(\K)\).  Let \(\Amr\subset\A\) be the set of
\(d\)-planes on which the restricted form has maximal possible rank.

\begin{theorem}
The maximal \(G\)-orbits in \(\operatorname{Ind}_P^G(\OO_0)\) are parametrized
by the isometry classes, inside \(\K\), of subspaces \(A\in\Amr\).
For the orbit \(\OO_A=G\cdot x_A\) attached to \(A\),
\[
m(\OO_A,P)=
\frac{|\pi_0(Z_G(x_A))|}{|\pi_0(Z_G(x_A)\cap P)|}
=
\begin{cases}
1,& \epsilon=+1,\\
1,& \epsilon=-1\text{ and }d\text{ is even},\\
2,& \epsilon=-1\text{ and }d\text{ is odd}.
\end{cases}
\]
Here \(\pi_0\) means connected components for the usual real topology.
\end{theorem}

\begin{proof}
First fix a non-degenerate lift \(V_1\subset\ker X_0\) of \(\K\), so that
\(\ker X_0=V_1\oplus R\) and \(V_1\simeq\K\).  Every element of
\(X_0+\uP\) is uniquely
\[
X_{C,B}=
\begin{pmatrix}
0&C^\vee&B\\
0&X_0&C\\
0&0&0
\end{pmatrix},
\qquad C:E'\to V_0,\quad B:E'\to E,
\]
where
\[
\lambda(C^\vee v,e')=-\langle v,Ce'\rangle_0,\qquad
\lambda(Bu,v)+\epsilon\lambda(Bv,u)=0.
\]
For \(D:E'\to V_0\), the element
\[
u_D=
\begin{pmatrix}
1&D^\vee&\frac12D^\vee D\\
0&1&D\\
0&0&1
\end{pmatrix}\in P
\]
satisfies
\[
u_DX_{C,B}u_D^{-1}
=X_{C-X_0D,\ B+D^\vee C-C^\vee D-D^\vee X_0D}.
\tag{1}
\]

For \(A\in\operatorname{Gr}_d(\K)\), let \(W_A\subset V_0/\im X_0\) be the
annihilator of \(A\) under
\((V_0/\im X_0)\times\ker X_0\to\mathbb R\).  Then \(\dim W_A=\kappa\).
Choose a lift \(V_A\subset V_0\) mapping isomorphically to \(W_A\), choose
\(S_A:E'\xrightarrow{\sim}V_A\), and put
\[
x_A:=X_{S_A,0}.
\]
The defining annihilator property gives
\[
\ker(S_A^\vee|_{\ker X_0})=A,
\qquad S_A^\vee|_{\ker X_0}:\ker X_0\to E\text{ is surjective}.
\tag{2}
\]

Conversely, take \(x=X_{C,B}\) and set
\(\phi_C=C^\vee|_{\ker X_0}\).  On the generic set where \(\phi_C\) is
surjective and \(L:=\ker\phi_C\) satisfies \(L\cap R=0\), the quotient
\[
A:=q(L)\subset\K
\]
has dimension \(d\).  A Levi conjugation first makes
\(\phi_C=S_A^\vee|_{\ker X_0}\).  Then \(C-S_A\) has image in
\(\im X_0\), so (1) makes \(C=S_A\).  A second conjugation by some
\(u_D\) with \(D(E')\subset\ker X_0\) kills the remaining \(B\)-block.
Thus every generic slice point is \(P\)-conjugate to \(x_A\).  The same
argument shows that \(x_A\) is independent, up to \(P\)-conjugacy, of the
choices of \(V_A,S_A\).  If two subspaces \(A,A'\subset\K\) are related by an
ambient isometry of \(\K\), that isometry lifts to an element of
\(Z_{G_0}(X_0)\), hence \(x_A\) and \(x_{A'}\) are \(P\)-conjugate.

The map \(C\mapsto \phi_C\) from \(\Hom(E',V_0)\) to
\(\Hom(\ker X_0,E)\) is surjective: modulo \(\im X_0\), it is exactly the
identification
\[
\Hom(E',V_0/\im X_0)\simeq \Hom(\ker X_0,E)
\]
coming from the two perfect pairings.  Therefore surjectivity of \(\phi_C\),
injectivity of \(\phi_C|_R\), and the condition \(q(\ker\phi_C)\in\Amr\) define
a dense subset \(\Omega\subset X_0+\uP\).  Moreover \(\Omega\) is a finite
disjoint union
\[
\Omega=\bigsqcup_i P\cdot x_{A_i},
\tag{3}
\]
where \(A_i\) run over the isometry classes in \(\Amr\).  These classes are:
non-degenerate symmetric \(d\)-planes classified by signature if
\(\epsilon=+1\); one non-degenerate symplectic class if \(\epsilon=-1\) and
\(d\) is even; and one alternating rank \(d-1\) class with a one-dimensional
radical if \(\epsilon=-1\) and \(d\) is odd.  The ambient transitivity in each
fixed class is Witt extension in the symmetric case and extension of symplectic
bases in the alternating cases.

Set \(Y=G\cdot(X_0+\uP)\).  For \(x_i=x_{A_i}\), the piece
\(P\cdot x_i\) is open in the affine slice, hence
\[
T_{x_i}(X_0+\uP)=T_{x_i}(P\cdot x_i)=[\mathfrak p,x_i]=\uP.
\]
For the action map \(G\times(X_0+\uP)\to\g\), the differential at
\((1,x_i)\) has image
\[
[\g,x_i]+\uP=[\g,x_i],
\]
so \(G\cdot x_i\) is open in \(Y\).  Since \(\Omega\) is dense in the slice,
\[
\operatorname{Ind}_P^G(\OO_0)=\overline{Y}
=\overline{G\cdot\Omega}
=\bigcup_i \overline{G\cdot x_i}.
\tag{4}
\]
Thus every orbit in the induced set lies below one of the listed orbits.
Because each \(G\cdot x_i\) is open in \(Y\), it has the dimension of
\(\overline Y\).  If another orbit had \(G\cdot x_i\) in its closure, the usual
boundary dimension inequality for locally closed real algebraic orbits would
force equality of the two orbits.  Hence the \(G\cdot x_i\) are exactly the
maximal induced orbits.

They are distinct precisely as stated.  For the normal form one computes
\[
\ker x_A=E\oplus A,\qquad
\im x_A=E\oplus\im X_0\oplus S_A(E'),\qquad
\ker x_A\cap\im x_A=E\oplus\rad(A).
\tag{5}
\]
Therefore
\[
\ker x_A/(\ker x_A\cap\im x_A)\simeq A/\rad(A)
\tag{6}
\]
as bilinear spaces.  This invariant recovers the signature in the symmetric
case and gives no further splitting in the two alternating cases, exactly as
the classification above says.

It remains to compute the multiplicity.  If \(A\) is non-degenerate, then by
(5)
\[
\ker x_A=E\oplus A,\qquad \ker x_A\cap\im x_A=E.
\]
Since \(E\) is the radical of the bilinear space \(\ker x_A\), every element of
\(Z_G(x_A)\) preserves \(E\).  Hence \(Z_G(x_A)\subset P\), and
\[
m(G\cdot x_A,P)=1.
\]
This covers \(\epsilon=+1\) and \(\epsilon=-1\) with \(d\) even.

Now assume \(\epsilon=-1\) and \(d\) is odd.  Write
\[
A=A_0\oplus \mathbb R r,\qquad \mathbb R r=\rad(A),
\]
with \(A_0\) symplectic, choose \(s\in\K\) with
\(\langle r,s\rangle_K=1\), and choose \(e\in E,e'\in E'\) with
\(\lambda(e,e')=1\).  Replacing \(x_A\) by a \(P\)-conjugate representative, we
may arrange that
\[
U=\mathbb Re\oplus\mathbb Rr\oplus\mathbb Rs\oplus\mathbb Re'
\]
is stable and
\[
x_A(e)=x_A(r)=0,\qquad x_A(s)=e,\qquad x_A(e')=r.
\tag{7}
\]
On \(W=U^\perp\), the kernel quotient is \(A_0\), so the non-degenerate case
gives \(Z_{G(W)}(x_A|_W)\subset P(W)\).

In the ordered basis \(e,r,s,e'\),
\[
x_A|_U=
\begin{pmatrix}0&I_2\\0&0\end{pmatrix},\qquad
J_U=
\begin{pmatrix}0&S_2\\-S_2&0\end{pmatrix},\qquad
S_2=\begin{pmatrix}0&1\\1&0\end{pmatrix}.
\]
Writing \(g=\begin{psmallmatrix}A&B\\ C&D\end{psmallmatrix}\), the equations
\(gx_A=x_Ag\) and \(g^TJ_Ug=J_U\) give
\[
C=0,\qquad D=A,\qquad A^TS_2A=S_2,\qquad
M^TS_2=S_2M\quad(M=A^{-1}B).
\]
Thus projection to \(A\) identifies \(\pi_0(Z_{G(U)}(x_A|_U))\) with
\(\pi_0(O(S_2))\).  Over \(\mathbb R\),
\[
O(S_2)=
\left\{
\begin{pmatrix}a&0\\0&a^{-1}\end{pmatrix}:a\in\mathbb R^\times
\right\}
\sqcup
\left\{
\begin{pmatrix}0&a\\a^{-1}&0\end{pmatrix}:a\in\mathbb R^\times
\right\}.
\]
The two diagonal branches and the two off-diagonal branches split further by
\(\operatorname{sgn}(a)\), so \(|\pi_0(O(S_2))|=4\).  Intersecting with
\(P(U)\), i.e. imposing preservation of \(\mathbb Re\), leaves only the
diagonal branch, hence two real components.  Therefore
\[
|\pi_0(Z_{G(U)}(x_A|_U))|=4,\qquad
|\pi_0(Z_{G(U)}(x_A|_U)\cap P(U))|=2.
\tag{8}
\]

Finally, relative to \(V=U\oplus W\), block projection from \(Z_G(x_A)\) to the
product of the \(U\)- and \(W\)-centralizers, and the analogous projection after
intersecting with \(P\), are surjective with connected fibers.  The same
\(W\)-factor occurs in numerator and denominator because
\(Z_{G(W)}(x_A|_W)\subset P(W)\).  By (8),
\[
m(G\cdot x_A,P)=4/2=2.
\]
This completes the proof.
\end{proof}

\end{document}
